\documentclass[10pt]{article}
\usepackage[margin=20mm]{geometry}
\usepackage{amsmath,amssymb,mathtools,fontspec,unicode-math,xurl,hyperref}
\IfFontExistsTF{DejaVu Serif}{\setmainfont{DejaVu Serif}}{\setmainfont{Latin Modern Roman}}
\setmathfont{latinmodern-math.otf}
\setlength{\parindent}{0pt}\setlength{\parskip}{6pt}\sloppy
\title{Planar Euler leapfrogging:
verification of the assigned author manuscript}
\author{AI-assisted technical review}
\date{}
\begin{document}
\maketitle
Abstract: This report examines arXiv:2311.15765v2, dated 5 December 2023. A principal-order Hilbert-conjugation error changes the degenerate Fourier block; correcting it reveals a periodic translation mode for every orbit parameter. The displayed approximate solution also retains a fourth-order residual, and the Nash–Moser induction does not justify its required uniform high norm. Each issue is tied to the relevant existence-proof input, with explicit calculations and counterevidence. Positive checks of point dynamics, geometric constraints and measure estimates are recorded separately. The conclusions concern this exact author version and do not assert that the nonlinear existence theorem is false or that the 2025 journal version contains the same defects.\par
\section*{Identity, outcome and scope}
Paper. Zineb Hassainia, Taoufik Hmidi and Nader Masmoudi, Rigorous derivation of the leapfrogging motion for planar Euler equations, arXiv:2311.15765v2, 5 December 2023 [V2]. The assigned PDF has 81 pages; all page references below use its printed page numbers, equal to PDF page numbers. The source SHA256 is recorded in Appendix 9 (\#app:evidence). The journal association is DOI 10.1007/s00222-025-01368-3, Inventiones mathematicae 242 (2025), 725–825. That 101-page publication is a different version. Its public metadata and subscription preview were checked, but a complete journal text was not compared with the assigned manuscript.\par
The full assigned text, including Appendix A and references, was read. The principal review checks cover the point-vortex Hamiltonian and period, the contour equation and leading derivative, both coordinate changes, the degenerate monodromy, the tame inverse, the induction exponents and the final parameter-measure argument. Actual source images resolve signs and coefficients that plain-text extraction obscures. Bounded native checks were integrated by the report writer; their scopes, numerical limitations and dependency versions are preserved in the evidence records.\par
Problems 1–3 establish distinct invalid steps in the central proof. The first is structural: a corrected leading periodic operator has a kernel throughout the parameter interval, so the displayed discrete-exception argument cannot supply its inverse. The second admits an explicit correction to the constructed approximation. The third admits a sufficient strengthened regularity choice and revised loss bookkeeping, conditional on a valid inverse. A separate stability-lemma issue is assessed below with its application limits. These are judgments of the supplied arguments, not a claim of nonexistence of leapfrogging patches. No author adoption is claimed: `scientific\_decision=pending`; `human\_reviews=[]`. The accompanying verification files record the independent focused checks and their completion status.\par
\section*{Quantifiers, spaces and the proof dependence}
Fix \(y_0>0\) and a compact interval \[O=[\xi_*,\xi^*],\qquad
  0<\xi_*<\xi^*<y_0/\sqrt2.\] Theorem 1.1 (p. 4) asserts that for every sufficiently small \(\varepsilon>0\) there is \(C_\varepsilon\subset O\) with \(|C_\varepsilon|\to |O|\), and for every \(\xi_0\in C_\varepsilon\) there is a four-patch Euler solution for all real time. The two upper patches have vorticity \(\varepsilon^{-2}\); their complex reflections have the opposite vorticity. Writing \(D_{t,k}^\varepsilon=z_k(t)+\varepsilon O_{t,k}^\varepsilon\), the normalized domains satisfy \[O_{t,2}^\varepsilon=O_{t+T(\xi_0)/2,1}^\varepsilon,\qquad
 O_{t+T(\xi_0),1}^\varepsilon=O_{t,1}^\varepsilon.\] The centers follow the translating point-vortex configuration. The theorem does not assert that the physical patches are periodic in a fixed laboratory frame, or include the collision and separatrix endpoints.\par
Use \(\varphi=\omega_0t\), \(\omega_0=2\pi/T\), and \(z_1-z_2=\sqrt q\,e^{i\Theta}\). The rotating radial parametrization is \[\gamma_1(\varphi,\theta)
  =e^{i\Theta(\varphi)}R(\varphi,\theta)e^{i\theta},
 \qquad R=\sqrt{1+2\varepsilon r}.\] The space \(H_0^s(\mathbb T^2)\) imposes zero \(\theta\)-average for each \(\varphi\); it imposes no centroid or first-harmonic constraint. The delay in the second upper patch is a shift by \(\pi\) in \(\varphi\), with the rotation sign treated in (3.12)–(3.18). The Hilbert convention on p. 23 is \[\mathcal He^{ij\theta}=i\operatorname{sign}(j)e^{ij\theta},\qquad j\ne0.\] All circle integrals below use normalized Haar measure. The notation \(a\leq_C b\) denotes \(a\leq Cb\) for a positive constant uniform in the small parameters under the stated hypotheses.\par
The parameter norm is \[\left\lVert h\right\rVert_s^{\operatorname{Lip}(\lambda)}
  =\sup_{\xi\in O}\left\lVert h(\xi)\right\rVert_{H^s}
   +\lambda\sup_{\xi\ne\eta}
       \frac{\left\lVert h(\xi)-h(\eta)\right\rVert_{H^{s-1}}}{|\xi-\eta|}.\] A scalar weighted Lipschitz bound therefore controls the ordinary Lipschitz constant only after division by \(\lambda\). This distinction matters in the measure estimate.\par
Lemma 3.3 supplies \(r_\varepsilon=g+O(\varepsilon)\) and claims \(\mathcal G(\varepsilon r_\varepsilon)=O(\varepsilon^5)\). The rescaled equation is \[\mathcal F(\rho)=\varepsilon^{-2-\mu}
       \mathcal G(\varepsilon r_\varepsilon+\varepsilon^{1+\mu}\rho)=0,\quad 0<\mu<1.\] Its derivative has a leading transport part, \(-\mathcal H/2\), and an order-\(\varepsilon^2\) operator supported on spatial modes \(\pm1\). Propositions 4.1–4.3 straighten the transport and identify that finite-rank operator; Proposition 5.1 inverts its degenerate block; Propositions 5.4–5.5 perturb and conjugate the approximate inverse. Section 6 then uses \[N_n=N_0^{(3/2)^n},\quad N_0=\varepsilon^{-\delta},\quad
 \lambda=\varepsilon^{2+\delta},\quad b_0=3-\mu .\] Thus a mistaken order-\(\varepsilon^2\) coefficient, an insufficient initial residual, or a missing small-norm input lies on a necessary link of the stated construction.\par
\section*{Principal-order conjugation and the periodic kernel}
Problem 1 (Major: the Hilbert conjugation changes the degenerate block).  The first identity on p. 40 and (4.31), propagated into Proposition 4.3, (5.4) and the monodromy analysis, use an incorrect factor and sign for the density pullback.\par
\subsection*{Exact identity and an independent Fourier check}
Equation reference: eq:exactH. For an orientation-preserving circle diffeomorphism \(f(\theta)=\theta+b(\theta)\), put \(p=f^{-1}\) and \(\mathcal Bh=f'(h\circ f)\), exactly as in (4.1). The stated Hilbert convention gives \[\mathcal Hh(\theta)
  =-2\partial_\theta\int_\mathbb T
     \log|e^{i\theta}-e^{i\eta}|\,h(\eta)\,d\eta.\] Changing variables in the integral and applying the inverse density pullback yields the exact formula \[
 (\mathcal B^{-1}\mathcal H\mathcal B-\mathcal H)h
  =-2\partial_\theta\int_\mathbb T
   \log\left|
    \frac{e^{ip(\theta)}-e^{ip(\eta)}}
         {e^{i\theta}-e^{i\eta}}\right|h(\eta)\,d\eta.\] The manuscript has \(+1\) in front of the same derivative and logarithm. The discrepancy is not a normalization ambiguity: both its explicit Fourier multiplier and its singular-kernel definition fix the sign and factor.\par
There is a direct smooth, real, zero-mean test. Take \(b=t\sin(2\theta)\), \(|t|<1/2\), so that \(f'>0\). Then \(\mathcal B_t=\operatorname{Id}+tD+O(t^2)\), with \(Dh=\partial_\theta(\sin(2\theta)h)\). For \(h=\cos\theta\), \[Dh=\tfrac12(\cos\theta+3\cos3\theta),\qquad
 \mathcal HDh=-\tfrac12(\sin\theta+3\sin3\theta),\] \[D\mathcal Hh=\tfrac12(\sin\theta-3\sin3\theta).\] Consequently \[\left.\frac{d}{dt}\right|_{t=0}
    (\mathcal B_t^{-1}\mathcal H\mathcal B_t-\mathcal H)\cos\theta=-\sin\theta.\] The printed formula instead gives \(+\frac12\sin\theta\). Equivalently, its coefficient on \(e^{i\theta}\) is \(+\frac{i}{2}e^{-i\theta}\) in place of \(-ie^{-i\theta}\). This witness satisfies the actual spatial-mean and smoothness conditions, and the change is a genuine small diffeomorphism.\par
\subsection*{Propagation in the actual orbit-dependent change}
Let \[\widehat\Theta=\Theta-\varphi,\quad
 a_2=\frac1q-\frac{e^{2i\Theta}}{w_3^2}
                  -\frac{e^{2i\Theta}}{w_4^2},\quad
 A=a_2e^{-2i\widehat\Theta}.\] Here \(w_3=i\sqrt q\sin\Theta+iy_0\) and \(w_4=\sqrt q\cos\Theta+iy_0\). The second harmonic in the actual first transformation is \(\varepsilon^2\Im(Ae^{2i\theta})\), after removing its spatially constant rotation. The manuscript’s own logarithmic expansion is \[K_0=-\frac{\varepsilon^2}{2}
        \Re\{A(e^{i\theta}+e^{i\eta})^2\}
          +O(\varepsilon^{2+\mu}).\] Write \(\widehat H_0h=\int_\mathbb T\Re(Ae^{i(\theta+\eta)})h(\eta)\,d\eta\). Zero mean removes one pure term and differentiation removes the other. Formula [eq:exactH] (\#eq:exactH) gives \[\mathcal B^{-1}\mathcal H\mathcal B-\mathcal H
       =2\varepsilon^2\partial_\theta\widehat H_0
            +O(\varepsilon^{2+\mu}),\] whereas (4.31) gives \(-\varepsilon^2\partial_\theta\widehat H_0\). Therefore the finite-rank operator in the conjugated linearization must be \[\mathcal Q_{\rm correct}=\mathcal Q_{\rm rot}+\widehat H_0,
 \qquad\hbox{rather than}\qquad
 \mathcal Q_{\rm printed}=\mathcal Q_{\rm rot}-\tfrac12\widehat H_0.\] This changes the principal \(\varepsilon^2\) block. At \(\varphi=0\), \(\Theta=\pi/2\), \(q=\xi_0^2\), \[a_2(0)=\xi_0^{-2}-(y_0+\xi_0)^{-2}-y_0^{-2}>0
       \quad (0<\xi_0<y_0/\sqrt2);\] the disputed coefficient is nonzero on the relevant orbits. Under the later scaling, \(\varepsilon^{2+\mu}+\varepsilon^6/\lambda=o(\varepsilon^2)\). It cannot be absorbed into the stated remainder.\par
\subsection*{The corrected block has a periodic translation mode}
Set \[S_w=w_3^{-2}+w_4^{-2},\qquad
 d=e^{-2i\widehat\Theta}/q,\qquad
 A=d-e^{2i\varphi}S_w.\] For the real function \(h=\cos(\varphi+\theta)\), its Fourier coefficients are \(h_1=e^{i\varphi}/2\), \(h_{-1}=e^{-i\varphi}/2\), and each changes sign under \(\varphi\mapsto\varphi+\pi\). Combining the delayed and undelayed kernels gives for the positive-frequency amplitude of \(\mathcal Q_{\rm correct}h\), after removing a common factor of \(1/2\), \[S_w h_1+(A-d)h_{-1}
  =S_w\frac{e^{i\varphi}}2
      -e^{2i\varphi}S_w\frac{e^{-i\varphi}}2=0.\] Also \((\partial_\varphi-\partial_\theta)h=0\). Thus the corrected order-\(\varepsilon^2\) operator on the degenerate subspace annihilates this nonzero, \(2\pi\)-periodic, real, spatial-zero-mean function for every \(\xi_0\in O\).\par
This is consistent with common horizontal translation of the four patches. Such a translation preserves the reflected-pair geometry and half-period relation. Neither \(H_0^s\) nor \(X^s\) removes it. An invertible periodic coordinate change cannot remove its Floquet multiplier \(1\). For the corrected unperturbed periodic system, \(\det(M_0(2\pi)-\operatorname{Id})=0\) throughout the interval; discarding a discrete set as in Lemma 5.2 cannot remedy this.\par
\subsection*{Counterevidence, impact and required work}
The calculation \(\det(e^{2\pi A_0}-\operatorname{Id})\simeq0.12126170473544845\) is reproducible for the printed constant matrix, with frequencies \(\sqrt{15}/4\) and \(\sqrt7/4\). It is not evidence for the corrected block. The corrected limiting characteristic polynomial is \(z^2(z^2+1)\). Independent orbit integration at five admissible parameters corroborates the translation mode, but the exact kernel calculation, not a floating-point near-zero determinant, establishes the issue.\par
The primary Hilbert-conjugation source [BM] uses a pure composition operator and the opposite Hilbert sign. Its smoothing statement cannot supply the manuscript’s density coefficient without rederivation. The explicit derivative check above does that rederivation.\par
This invalidates the displayed nonresonance route from Proposition 4.3 to the approximate inverse used in Section 6. It does not disprove nonlinear existence. A repair would require a correctly constrained phase space and the associated solvability conditions, or another treatment of the neutral mode, followed by a new inverse and iteration analysis. No such complete repair is claimed here. The kernel and conjugation error are one dependent issue, not two separately counted Major findings.\par
\section*{The residual of the displayed approximate solution}
Problem 2 (Major: the stated construction leaves an order-\(\varepsilon^4\) term).  In Lemma 3.3, equations (3.49)–(3.53), the definition of \(A_1\) at the top of p. 32 omits a term needed for the asserted \(O(\varepsilon^5)\) residual and the initialization (6.5).\par
Equation reference: eq:A1. Put \(D_0=(\partial_\theta-\mathcal H)/2\). The construction chooses \(r_0=g+\varepsilon B_0\), with no first harmonic, so that the circle-state forcing is canceled. Expansion using the correctly printed final derivative (3.21) gives the order-\(\varepsilon^4\) coefficient \[
 A_{1,\rm correct}
  =\omega_0\partial_\varphi g
     -\omega_0\dot\Theta\,\partial_\theta g
     -\tfrac34\partial_\theta(g^2).\] The middle term is missing from \(A_{1,\rm printed}\). The \(B_1\) on p. 32 contains the rotation coefficient multiplying \(B_0\), which contributes at the next order; it does not contain the missing term multiplying \(g\).\par
For a completely admissible parameter test take \(y_0=1\), \(\xi_0=1/2\) in, for example, \(O=[0.45,0.55]\). At \(\varphi=0\), \[q=\tfrac14,\quad \Theta=\tfrac\pi2,\quad
 a_2=\tfrac{23}{9},\quad
 \omega_0\dot\Theta=4+\tfrac43=\tfrac{16}{3}.\] The last equality follows directly from the point-vortex ODE; it does not approximate the frequency. At \(\theta=\pi/4\), \(g=(23/9)\cos2\theta\), so \[-\omega_0\dot\Theta\,\partial_\theta g
       =\frac{736}{27}\ne0.\] The construction printed in the proof sets \(D_0r_1=-\varepsilon A_{1,\rm printed}\) and \(r_\varepsilon=r_0+\varepsilon r_1\). It therefore leaves \[\mathcal G(\varepsilon r_\varepsilon)(0,\pi/4)
       =\frac{736}{27}\varepsilon^4+O(\varepsilon^5),\] in place of the claimed fifth-order bound.\par
\subsection*{An explicit verified local correction}
All terms in [eq:A1] (\#eq:A1) have even spatial modes. On the second harmonics, \(D_0g=\partial_\theta g/4\). Replacing the printed \(r_1\) by \[r_{1,\rm corrected}
      =r_{1,\rm printed}+4\varepsilon\omega_0\dot\Theta\,g\] gives \(D_0r_{1,\rm corrected}=-\varepsilon A_{1,\rm correct}\). Equivalently, add \(4\varepsilon^2\omega_0\dot\Theta\,g\) to \(r_\varepsilon\). This preserves its zero mean, absence of first harmonics, smooth parameter dependence and stated size; substituting it into the Taylor expansion cancels exactly the missing fourth-order term. The unchanged higher-order bounds then give the asserted fifth-order residual for this corrected approximation.\par
\subsection*{Why the error is material despite this correction}
Without the correction, \(\mathcal F(0)\) has order \(\varepsilon^{2-\mu}\), not \(\varepsilon^{3-\mu}\) as in (6.5). Multiplication by the stated inverse scale \(\lambda^{-1}=\varepsilon^{-2-\delta}\) gives \(\varepsilon^{-\mu-\delta}\), instead of the small initial size \(\varepsilon^{1-\mu-\delta}\). Thus the printed object does not supply the small-data input advertised for the core Nash–Moser construction. The finding concerns that invalid construction, not falsity of the existential assertion of Lemma 3.3. The explicit correction is counterevidence against a broader nonexistence claim; its simplicity does not erase the demonstrated dependence on the incorrect residual. It also does not fix Problem 1.\par
\section*{Closing the high norm in the induction}
Problem 3 (Major: the inverse’s high-norm hypothesis is not obtained).  The interpolation and summation on p. 56 do not establish the uniform norm required by (5.49). The parameter restrictions (6.4) force the displayed interpolated exponent to be positive. The subsequent use of a smaller derivative loss in (6.31) also requires replacement.\par
Let \[L=2s_0+2\tau+3,\qquad
 R_*=2s_0+2\tau+4+\tfrac32\mu_2,\qquad
 A_\varepsilon=\varepsilon^{b_0}/\lambda=\varepsilon^{1-\mu-\delta}.\] The induction supplies increments of size \(A_\varepsilon N_{k-1}^{-a_2}\) in \(H^L\), while the \(H^{b_1}\) norm of the iterates is permitted to grow as \(A_\varepsilon N_{n-1}^{\mu_1}\). The inverse needs \(\left\lVert \rho_n\right\rVert_{R_*}^{\operatorname{Lip}(\lambda)}\le1\). Page 56 uses half-weight interpolation and obtains \[\left\lVert u_k\right\rVert_{R_*}^{\operatorname{Lip}(\lambda)}
   \mathrel{\leq_C}A_\varepsilon N_{k-1}^{(\mu_1-a_2)/2}.\] It then sums this upper bound as though it decayed. However, (6.4) gives \[\mu_1>\frac6\delta>
       \frac{6(a_1+2)}{1-\mu}>6a_1>9a_2 .\] The exponent is strictly positive for every permitted choice.\par
\subsection*{An explicit test of the implication used on p. 56}
Take \[(s_0,\tau,\mu,a_2,a_1,\mu_2,\delta,\mu_1,b_1)
   =(4,2,\tfrac14,4,60,41,\tfrac1{100},630,1400).\] Every strict inequality in (6.4) holds: \(4>3\), \(60>57\), \(41>40\), \[\tfrac1{100}
 <\min\{\tfrac14,\tfrac3{248},\tfrac3{56}\},\qquad
 630>627,\qquad 1400>\max\{1315,149\}.\] Here \(L=15\), \(R_*=155/2\), and the printed exponent is \(313\). The \(4\mu_1/3\) term in the source gives the lower bound \(1315\). The earlier retained witness with \(b_1=900\) was therefore inadmissible and is superseded by the present parameters; its historical computation is not evidence of an admissible example.\par
For fixed small \(\varepsilon\), let \(m_k=\lfloor N_{k-1}^{1/4}\rfloor\), ignoring finitely many initial repetitions if needed, and let \(w_m=\cos(m\theta)\) normalized to have \(H^{15}\) norm one. Set \[u_1=0,\qquad
 u_k=A_\varepsilon N_{k-1}^{-4}w_{m_k}\ (k\ge2),\qquad
 \rho_n=\sum_{k=2}^n u_k.\] These are real, spatial-zero-mean trigonometric polynomials. They are independent of the parameter and of \(\varphi\), so their Lipschitz seminorm vanishes. Their frequency support satisfies \(\rho_n\in E_{n-1}\). Their low norms are summable, and \[\left\lVert u_k\right\rVert_{1400}\asymp
       A_\varepsilon N_{k-1}^{1369/4},\qquad
 \left\lVert u_k\right\rVert_{155/2}\asymp
       A_\varepsilon N_{k-1}^{93/8}.\] The superexponentially spaced high-norm sum is bounded by a constant times \(A_\varepsilon N_{n-1}^{1369/4}\le A_\varepsilon N_{n-1}^{630}\). Thus the low/high bounds and localization used in that interpolation argument hold. But distinct Fourier modes are orthogonal, and \(\left\lVert \rho_n\right\rVert_{155/2}\) is unbounded.\par
This is a counterexample to the displayed implication from those bounds, not an asserted Euler solution or a sequence satisfying the separate nonlinear residual condition (P2). The page-56 argument does not invoke (P2) to obtain its high-norm conclusion. Any alternative use of the nonlinear structure must actually be supplied.\par
\subsection*{Sufficient repair conditions and the genuine losses}
Equation reference: eq:b1. Since the authors can choose \(b_1\), a valid interpolation can be arranged by strengthening its lower bound. Set \[\vartheta=\frac{R_*-L}{b_1-L},\qquad
 \kappa=a_2-\vartheta(\mu_1+a_2).\] Interpolation of the actual low increment bound and the high increment bound gives \[\left\lVert u_k\right\rVert_{R_*}^{\operatorname{Lip}(\lambda)}
       \mathrel{\leq_C}A_\varepsilon N_{k-1}^{-\kappa}.\] Both the supremum norm and the one-derivative-lower difference norm interpolate with the same \(\vartheta\). The sufficient condition \[
 b_1>L+2\frac{\mu_1+a_2}{a_2}(R_*-L)\] ensures \(\kappa>a_2/2>0\), so the increments sum in \(H^{R_*}\). For the same remaining parameters, \(b_1=20000\) gives \(\vartheta=25/7994\) and \(\kappa=8063/3997>2\). The available upper regularity \(S\) must be selected to include this \(b_1\) and the finite losses, for example \(S\ge b_1+4\tau+6\) for the displayed bounds used here. No fixed-\(\varepsilon\) infinite-smoothness conclusion is inferred from this finite-regularity argument.\par
The actual inverse estimate (6.6) is \[\left\lVert T_nh\right\rVert_s\mathrel{\leq_C}\lambda^{-1}
 \bigl(\left\lVert h\right\rVert_{s+2\tau}
       +\left\lVert \rho_n\right\rVert_{s+4\tau+5}\left\lVert h\right\rVert_{s_0+2\tau}\bigr).\] Equation (6.31) replaces \(s+4\tau+5\) by \(s+2\tau\); that general replacement is unjustified. After [eq:b1] (\#eq:b1), low-order use can instead exploit \[R_*-(L+4\tau+5)=\tfrac32\mu_2-4\tau-4>0.\] For \(s\le L\) this gives \[\left\lVert u_{n+1}\right\rVert_s
   \mathrel{\leq_C}A_\varepsilon
          N_n^{s-s_0+2\tau-\frac23a_1}.\] At \(s=b_1\), retain \(N_{n-1}^{4\tau+5}\left\lVert \rho_n\right\rVert_{b_1}\) and combine it with the residual decay \(N_{n-1}^{-a_1}\). The resulting two high-growth exponents are bounded by \[\tfrac23\mu_1+\tfrac43\tau+\tfrac73+\tfrac2\delta,
 \qquad
 \tfrac23\mu_1+2\tau+\tfrac23(4\tau+5-a_1),\] both strictly below \(\mu_1\) under (6.4). The original reviewer calculation record gives the complete inequality chain.\par
The apparent loss of \(2\tau\) in (6.22) has a bounded explanation: retaining the required \(H^{s_0+2\tau+2}\) norm yields a factor \(N_n^{-\frac32\mu_2+2\tau+2}\), which is at most \(N_n^{-\mu_2+2}\) because \(\mu_2\ge4\tau+3\). Its final weaker exponent can survive, although its intermediate display needs correction. The small powers \[1-\mu-\delta,\quad \mu-\delta,\quad
 2+\mu-\delta(1+\mu_2),\quad
 1-\mu-\delta(a_1+2)\] are all positive for the specified choices.\par
This supplies a precise route for the affected regularity and loss estimates, conditional on the nonlinear and inverse inputs. It does not make the printed half-weight summation valid or establish the claimed freedom to use every choice satisfying only (6.4). The uniform high norm is required before invoking the inverse and again when evaluating the limiting frequency coefficient. Its absence is therefore a core proof gap. Neither this conditional repair nor a larger \(b_1\) repairs the neutral-mode obstruction in Problem 1.\par
\section*{The normalized transport-stability statement}
Problem 4 (Major in the stated stability bridge; application not refuted).  Proposition A.9(4), p. 72, gives a state-difference estimate for the normalized correction \(c\) with the scaling of a full transport speed. With the proof’s definition \(c_0=0\), a constant-coefficient example contradicts that estimate. The bridge to (4.26) requires a corrected statement and explicit difference hypotheses.\par
Equation reference: eq:A9literal. Appendix A starts from \[T_0=\epsilon_1\Omega\,\partial_\varphi+
       \partial_\theta[(c_{-1}+\epsilon_2f_0)\,\cdot],\] and concludes with full speed \(c_{-1}+\epsilon_2c\). Thus \(c\) is the normalized correction. Part (4) asserts, uniformly in the small parameters, \[
 \left\lVert \Delta c\right\rVert^{\operatorname{Lip}(\lambda)}
 \mathrel{\leq_C}\left\lVert \Delta c_0\right\rVert^{\operatorname{Lip}(\lambda)}
       +\epsilon_2\left\lVert \Delta\rho\right\rVert_{2s_0+2\tau+3}^{\operatorname{Lip}(\lambda)} .\] The proof on the same page initializes \(c_0=0\).\par
Take \(\Omega=1\), \(c_{-1}=1/2\), and \(f_0(\rho)=\rho\) for constant states \(\rho=r\), with \(r\) small enough that all stated Sobolev norms are at most one. For instance, take \(r_1=1/4\), \(r_2=0\) in the paper’s normalized Fourier norm. The absolute bounds (A.10)–(A.11) hold with \(\sigma_0=0\). For \(\epsilon_1=t\), \(\lambda=t^2\), \(\epsilon_2=t^4\) and fixed sufficiently large \(N_0\), \[\epsilon_1^{-1}\lambda+
       N_0^{\mu_2}\epsilon_2\lambda^{-1}
        =t+N_0^{\mu_2}t^2\longrightarrow0 .\] All smallness conditions (A.13) therefore hold as \(t\to0\). No spatial-mean restriction on \(\rho\) is imposed in this general appendix proposition.\par
The homological forcing has no nonconstant Fourier modes. The recurrence (A.26) gives \(\beta=0\) and \(c=r\) exactly after the first mean update. Consequently [eq:A9literal] (\#eq:A9literal) would require \[|r_1-r_2|\mathrel{\leq_C}t^4|r_1-r_2|,\] which is impossible with a uniform constant. This is an exact test of the stated normalized lemma and its construction, not a numerical approximation to the patch solution.\par
Two plausible intended readings limit the conclusion. If \(c_0\) in part (4) meant \(\langle f_0\rangle\), the constant test would no longer contradict it; however, that is not the \(c_0=0\) defined by this proof and needs a new definition. If the left side meant the full speed \(c_{-1}+\epsilon_2c\), the factor \(\epsilon_2\) would have the expected scale. That reading changes the displayed quantity. Moreover, (A.10)–(A.11) give absolute state norms, not the state-difference assumptions needed for a general \(\Delta\rho\) conclusion.\par
The source invoked for the omitted proof, Hmidi–Roulley [HR21], Proposition 6.2 and (6.30), estimates differences of the full transport coefficient. Its parameter regularity and frequency cutoffs also differ from those in A.9. The actual primary statement was inspected; a citation to it does not resolve these normalization and hypothesis changes by itself.\par
In the specific application, (4.6)–(4.7) provide additional difference controls. A corrected estimate for the full coefficient could therefore still yield (4.26). The constant example does not disprove (4.26). The material issue is that the general lemma stated as its justification is false under its explicit normalization, and the necessary adapted stability proof is not supplied. Section 6.2 uses (4.26) to compare the limiting and finite-step Cantor inequalities, so this is a core dependency rather than an unrelated appendix assertion. A complete corrected stability lemma has not been verified here; its exact hypotheses, loss indices and parameter factors require restatement and proof.\par
\section*{Positive checks, counterevidence and local corrections}
\subsection*{Point dynamics and the compact parameter range}
For \(Y=y_0\), the reduced Hamiltonian is \[H(\eta,\xi)=\tfrac12\log
       \frac{(Y^2-\xi^2)(Y^2+\eta^2)}{\eta^2+\xi^2}.\] Direct differentiation reproduces (2.4) with \(\dot\eta=\partial_\xi H\), \(\dot\xi=-\partial_\eta H\). If \(\alpha=\xi_0^2/Y^2\), \(v=1-2\alpha\), \(x=\xi^2/Y^2\) and \(z=\eta^2/Y^2\), its level set gives \[x(1+\alpha z)=\alpha-vz,\qquad
 \frac{\eta^2+\xi^2}{Y^2}
       =\frac{\alpha(1+z)^2}{1+\alpha z}.\] Hence \(\xi_0^2\le\eta^2+\xi^2\le\xi_0^2/v\) and \(|\xi|\le\xi_0<Y/\sqrt2\). The regular compact reduced orbit is noncolliding and periodic; the centroid has bounded horizontal velocity. These facts establish complete relative-periodic point trajectories, not laboratory-frame boundedness.\par
For \[K(\alpha,u)=
 \frac{(1-\alpha)(v+\alpha u)^2}
      {v(v+\alpha^2u)^{3/2}},\qquad
 T=4\xi_0^2\int_0^1\frac{K(\alpha,s^2)}{\sqrt{1-s^2}}\,ds ,\] the denominators remain positive on compact parameter intervals. Dominated analytic differentiation justifies analyticity of \(T\) and \(\omega_0\). Writing \(A=v+\alpha u\), \(B=v+\alpha^2u\), the numerator of \((1-\alpha)vAB\,\partial_\alpha\log K\) is \[\alpha v^2(1-u)^2+
 2v(1-\alpha)(1-\alpha/2+\alpha^2)u(1-u)
 +2(1-\alpha)^4u^2>0.\] Also \[\partial_u\log K
   =\frac{\alpha[(4-3\alpha)v+\alpha^2u]}{2AB}>0 .\] It follows that \[2\pi\xi_0^2\le T\le\frac{2\pi\xi_0^2}{v},\qquad
 T'\ge4\pi\xi_0,\qquad
 |\omega_0'|\ge\frac{2v^2}{\xi_0^3}.\] This supplies the needed twist on every fixed compact \(O\). The record also derives the angular estimates independently.\par
The reviewer checked physical-time orbit integration, quadrature and the correctly normalized elliptic formula at \(\xi_0=0.02,0.1,0.3,0.5,0.65\), \(Y=1\). The largest relative orbit/quadrature period difference was \(2.253\times10^{-12}\). These finite checks corroborate the analytic argument; they are not interval bounds or a simulation of patch existence. Behring–Goodman [BG] was checked with its coordinate and time rescalings. Its neutral and transverse stability directions cannot be confused with an existence obstruction for the present symmetric ansatz.\par
The p. 13 elliptic parameter must be \(\alpha=\xi_0^2/Y^2\); the names of the two standard complete elliptic integrals are reversed there. The derivative bracket on p. 15 also needs the ratio \(A/B\) rather than the printed product. The positive derivative identities above verify that the period/twist conclusions survive these local formula corrections. They are not additional Major findings. The unavailable original Love PDF and the cited table entries were not represented as read; the required claims were derived independently.\par
\subsection*{Contour, geometry and non-rigidity}
Direct differentiation of the quadratic cross-interaction kernels agrees with the leading delayed and reflected terms in (3.24). The actual final formula (3.21) contains \(-\varepsilon^2g/2\); its coefficient is correct despite ambiguous text extraction. The intermediate (3.44) loses a half that the final formula restores. The delayed kernel in p. 34’s definition of \(K_{0,1}^\varepsilon\) should be \(K_3^\varepsilon\), without the extra \(1/\varepsilon\). These local displays are distinguished from the incorrect principal conjugation in Problem 1. The transport change itself can be checked by the chain rule, and its density factor preserves the spatial integral. A spatially constant, non-small rotation does not require the small-diffeomorphism lemma when separated from the small \(\theta\)-dependent part.\par
Conditional on a valid small solution, zero spatial mean gives \(|O_{t,k}^\varepsilon|=\pi\) exactly and hence circulations \(\pm\pi\). The point-center distances are uniformly bounded below by \[d_*=\min\{\xi_*,\,y_0-\xi^*\}>0.\] For \(R\le2\), \(4\varepsilon<d_*\) makes all four patches disjoint for all time. The same smallness ensures \(R^2>0\). The rescaled bound \(\rho=O(\varepsilon^{1-\mu-\delta})\), with \(\delta<1-\mu\), is consistent with the leading \(\varepsilon^2\) shape term and the stated \(O(\varepsilon^{2+\mu})\) remainder.\par
A separate exact check supports leading non-rigidity throughout the admissible interval. Put \(s=\xi_0/y_0\). At the two vertical-axis crossings, the second-harmonic coefficients, apart from \(y_0^{-2}\), are \[a_+=s^{-2}-(1+s)^{-2}-1,\qquad
 a_-=s^{-2}-(1-s)^{-2}-1.\] At a horizontal-axis crossing, set \(Q=s^2/(1-2s^2)\). The squared amplitude is \[H_*=\left[Q^{-1}+1-\frac{Q-1}{(Q+1)^2}\right]^2
                   +\frac{4Q}{(Q+1)^4}.\] The reduced numerator polynomials of \(a_+^2-a_-^2\) and \(H_*-a_+^2\) have greatest common divisor one; the saved exact calculation verifies a Bezout identity. All denominators are nonzero on \(0<s<1/\sqrt2\), so the three amplitudes cannot all agree at any admissible parameter. For example, at \(s=1/2\), \(a_+=23/9\), \(H_*=97/9\), and \(H_*-a_+^2=344/81\). This is a positive shape check conditional on construction, not a substitute for the missing inverse.\par
\subsection*{Normal modes, extension and parameter measure}
The cutoff inverse on \(|j|\ge2\) has the asserted \(\lambda^{-1}\) size with a \(2\tau\) spatial loss in the weighted Lipschitz norm. In (5.34), the term \(\Pi_{N_n}^{\perp}D_1T_n\) is zero by its Fourier support; the remaining high-frequency projector has the stated decay. An incorrect sign in the preceding decomposition does not change this support calculation. The actual higher coefficient losses are tracked in Problem 3.\par
Extension from a closed subset of a real parameter interval can be performed by affine interpolation across complementary gaps, and constant continuation at outside endpoints. The same linear construction preserves all relevant Hilbert norms and the Lipschitz bounds. Thus no incompatibility between separate Sobolev extensions needs to be presumed from the brief Kirszbraun citation.\par
The p. 67 deduction of an ordinary, parameter-uniform bound on \(c_2'\) from (5.6) drops a factor \(\lambda^{-1}\). A correct finite-difference argument uses \(\operatorname{Lip}(c_2)\mathrel{\leq_C}\lambda^{-1}\). In the resonant band \(|j|\le24a\varepsilon^2|\ell|\), with \(a=\omega_0(\xi_*)\) and a fixed compact twist lower bound \(c_{\rm tw}>0\), it gives \[|f_{\ell,j}(\xi)-f_{\ell,j}(\eta)|
 \ge |j|\left(\frac{c_{\rm tw}}{24a}
          -C\varepsilon^2-C\frac{\varepsilon^3}{\lambda}\right)|\xi-\eta|.\] Here \(\varepsilon^3/\lambda=\varepsilon^{1-\delta}\to0\). Thus the lower Lipschitz bound needed for the measure estimate does survive. Each resonant interval has length at most \(C\varepsilon^{2+\delta}|j|^{-1-\tau}\); summing the \(O(|j|\varepsilon^{-2})\) relevant time indices and using \(\tau>1\) gives \(O(\varepsilon^\delta)\). This closes that local loss issue, conditional on the frequency stability and solution inputs. It cannot restore the discrete-exception treatment of the corrected degenerate block.\par
\section*{Overall judgment and remaining work}
The main existence argument in this assigned version is not verified by the presented proof. Problem 1 requires a substantive reanalysis of neutral directions, range conditions and the approximate inverse. Problem 2 has a verified change of the approximate solution. Problem 3 has sufficient additional regularity conditions and conditional replacements for the affected induction estimates. Problem 4 identifies a false normalized stability statement and an incompletely specified application bridge; it does not establish that the desired application estimate is false. Local formula and norm corrections with valid replacements are not promoted to separate core failures.\par
The review does not claim the 2025 journal text is unchanged. The official version/correction lookup conducted for this review did not return a separate correction; that limited search is not evidence that no correction exists. A journal-version comparison and a revised neutral-mode proof remain outside the completed bounded checks. The current independent acceptance is recorded separately in the verification files. No empirical experiment or author-code execution is needed to validate the exact countercalculations above, and absence of code is not a defect in this mathematical paper.\par
Item | Judgment | Repair and remaining scope\par
Problem 1 | Major | Correct coefficient and periodic kernel verified; full neutral-mode repair not supplied.\par
Problem 2 | Major | Explicit even-mode correction verified; does not repair Problem 1.\par
Problem 3 | Major | Stronger \(b_1\) and genuine-loss estimates verified conditionally; inverse still unresolved.\par
Problem 4 | Major | Literal normalized lemma fails; actual (4.26) requires a corrected stability bridge.\par
Local formulas | Minor/local | Period, kernel-display and weighted-norm corrections have bounded replacements; no additional core failure asserted.\par
Scientific decision | Pending | No human review or author adoption; focused acceptance is recorded in the verification files.\par
\section*{Supporting checks}
The point-dynamics and translation-kernel records include exact identities, original calculation code, finite numerical residuals, source/dependency hashes and actual image inspections. The linearization record contains the precise density conjugation, an independent Fourier computation, the approximate-solution correction and the primary-source convention comparison. The historical iteration record retains the original exponent ledger and stability/measure qualifications. The root record contains the full-text read ranges, independent physical geometry, the all-parameter quadrupole certificate and additional actual image observations. All retained evidence images are tied to the inspecting role and their bytes; rendering alone is not recorded as inspection.\par
The assigned PDF SHA256 is\par
`5c3ecff79051ab5bff9baebe2a634c00b5cd46f6045adf65337911aca9ad0afe`.\par
The assigned source-manifest SHA256 is\par
`52dd2a9698b551e1b69681010ab5a96323af3adbdfcdbec8b4848d5d64606a66`.\par
The full extracted-text SHA256 is\par
`b4178988895ae121a0b4bebefdcff2eefea40702ff6e364c3db975b24d47d126`.\par
Its five orbit cases have maximum period discrepancy \(2.253\times10^{-12}\), corrected-monodromy translation residual \(2.020\times10^{-11}\), and pointwise kernel residual \(1.364\times10^{-14}\). The independent exact induction calculations ran in process 54284, exit zero. The root’s physical calculation ran in process 2448, exit zero. These are reviewer-authored calculations, not author code.\par
The first root polynomial-certificate attempt, process 65236, failed an overly strong assertion because unreduced rational numerators retained the excluded endpoint factor \(s^2\). The diagnostic run 43048 identified that factor. The final run 43260 canceled the rational expressions, normalized their numerator polynomials and verified gcd one together with an exact Bezout identity; it exited zero. The failed script and diagnostic are retained and are not counted as successful certificates. The first report build, process 52512, exited one because the cached bundle lacked the physical font `cmex8.pfb`. Its source, logs and unchanged dependency hashes are retained. The next build, process 12484, resolved that extension-font issue but found the physical AMS inequality font unavailable. The final source uses the cached extension font at all requested sizes and the explicitly defined \(\leq_C\) notation for constant-factor inequalities; no font was installed. Other utility/retrieval failures are preserved in the bounded roles’ logs. An inaccessible public historical source was not represented as read.\par
The report is built from the current saved TeX using the existing authorized Tectonic 0.17.0 executable with `–untrusted –only-cached –keep-logs –keep-intermediates` and an explicit dependency file. No author code, shell escape, new installation or credentials are used. Their commands, paths and identities remain historical. The cache inventory is a hashed superset; `build.d` records direct dependencies. These post-build hashes are not fabricated inside this pre-build source.\par
\section*{Supporting checks}
The preserved recovery conversation’s root performed the full source-text reading, integrated the original bounded analyses, and built and inspected that recovery report. The present independent acceptance rechecks the four reported problems and their limits, corrects the inadmissible parameter witness, and separately records its actual new build and all-page inspection. Their prior roles are not misrepresented as new external sessions.\par
The present acceptance records its own controller UI observation separately in provenance. Neither observation establishes an unseen backend identity. No external Work/Codex or replacement-model session was opened. Native checks are not independent human acceptance.\par
the recovery token is\par
In that recovery, the packet and 82 materialized preserved files were fully parsed and hash-checked with their snapshot counterparts. The 440-entry snapshot index was fully parsed; this is not a claim that all 440 payloads were newly hashed. The archived old conversation has no invented final receipt. Those are historical recovery counts, not this acceptance’s full frozen-input check. The current acceptance writes only its exclusive staging; its actual input, disk, build and inspection records retain their own token and execution identity.\par
This is a systematic, bounded mathematical review rather than a formal proof-assistant verification of every external theorem. Core dependencies were checked in the exact consulted versions, and unavailable historical sources were replaced only where an independent derivation is supplied. The numerical checks are finite and nonvalidated. One borderline derivative economy in (4.35) was not independently optimized: applying the generic primary smoothing-kernel lemma to the printed change estimate naively costs one additional derivative. This is recorded as a limitation, not a demonstrated additional error or a silently verified estimate. Unread portions of the external references and reference [27] itself are not claimed checked. No whole-paper correctness, author adoption or human endorsement is claimed.\par
9\par
Z. Hassainia, T. Hmidi and N. Masmoudi, Rigorous derivation of the leapfrogging motion for planar Euler equations, arXiv:2311.15765v2, 5 December 2023, 81 pp. https://arxiv.org/abs/2311.15765v2. Exact reviewed source.\par
Z. Hassainia, T. Hmidi and N. Masmoudi, same title, Invent. Math. 242 (2025), 725–825, DOI:10.1007/s00222-025-01368-3. https://doi.org/10.1007/s00222-025-01368-3. Public metadata/preview checked; full-text equivalence not established.\par
M. Berti and R. Montalto, Quasi-periodic standing wave solutions of gravity-capillary water waves, Mem. Amer. Math. Soc. 263 (2020), no. 1273. Consulted institutional accepted manuscript, PDF pp. 49–51 (printed pp. 43–45), especially (2.106), (2.110) and Lemma 2.36. https://air.unimi.it/bitstream/2434/726528/2/Berti-Montalto-Memoires.pdf. The exact file hash is bound in the dependency ledger.\par
T. Hmidi and E. Roulley, Time quasi-periodic vortex patches for quasi-geostrophic shallow-water equations, arXiv:2110.13751v1. https://arxiv.org/abs/2110.13751v1. Consulted Proposition 6.2, its assumptions and relevant KAM construction; no equivalence to later versions assumed.\par
Z. Hassainia and E. Roulley, Boundary effects on the emergence of quasi-periodic solutions for Euler equations, arXiv:2202.10053v1. https://arxiv.org/abs/2202.10053v1. Consulted PDF pp. 24–27 for the difference-quotient and tame kernel input; no whole-source review claimed.\par
B. M. Behring and R. H. Goodman, Stability of leapfrogging vortex pairs: A semi-analytic approach, Phys. Rev. Fluids 4 (2019), 124703. https://doi.org/10.1103/PhysRevFluids.4.124703. Author-hosted publisher version consulted with explicit time and coordinate normalization; relevant PDF pp. 3, 5, 6 and 8 inspected.\par
\end{document}
